% chapters/03_quantitative_wellposedness.tex
%
% NOTE ON NUMBERING: in the original document this chapter is internally
% lettered "Appendix B" (Theorem B.9, Assumptions B.1-B.7, etc.), reflecting
% an appendix-style numbering convention used throughout the thesis for its
% main chapters. Since this is genuinely Chapter 3 in the thesis's actual
% sequence (as referenced everywhere else: Table 1.1, the abstract, the
% roadmap), the theorem/assumption environments below are renumbered to the
% chapter's own counters (Theorem 3.1, Assumption 3.1, etc.) via the shared
% preamble's [chapter] counter reset, with the original "B.x" label
% preserved in a footnote/remark at first use for cross-reference with any
% earlier excerpted material (e.g. Lean 4 listings) that still cite "B.9".

\section{Overview}

This chapter proves Theorem~\ref{thm:wellposedness} (originally labelled Theorem~B.9),
the central structural result of the thesis. We establish that the forward-backward
stochastic differential equation (FBSDE) system \eqref{eq:fbsde-system} characterising
the mean-field game with common noise admits a unique solution under the standing
assumptions of Section~1.9. Moreover, we provide an explicit quantitative stability
estimate: the solution map is a contraction in a weighted Wasserstein-2 metric on the
space of measure flows, with a contraction constant that depends explicitly on the
model parameters --- the Lipschitz constant of the drift $L_b$, the Frobenius norm of
the common-noise coefficient $\|\sigma_0\|_F$, and bounds on Lions derivatives. The
$\varepsilon$-Nash property for the finite-player game is also verified.

While the existence and uniqueness of solutions to common-noise MFGs have been
established under various conditions --- notably by \citep{ahujarenyang2019} under
convexity and weak-monotonicity, and in the comprehensive probabilistic treatment of
\citep{carmonadelarue2018} --- the present contribution is the derivation of explicit
parameter-dependent stability constants suitable for numerical analysis and
reinforcement learning algorithms. These constants quantify precisely how the model
parameters affect the contraction rate. \textbf{We emphasize that the contraction is
established in expectation, not pathwise, and relies on the synchronous coupling of two
candidate solutions sharing the identical common noise.}

The proof proceeds by formulating the MFG equilibrium as a fixed-point problem on the
space of measure flows, $\Phi(\mu) = \mu$, where $\Phi$ maps a candidate measure flow to
the conditional law of the optimally controlled state. Using the Stochastic Maximum
Principle (Section 2.7) and the theory of McKean--Vlasov FBSDEs, we show that for any
given measure flow $\mu$, there exists a unique optimal control and corresponding state
process. The key technical step is the derivation of a Wasserstein stability estimate
for the map $\Phi$ in expectation, which relies on synchronous coupling of two candidate
solutions, the Lions derivative calculus, and It\^o's formula applied to the squared
Wasserstein distance between the conditional laws. Under sufficient regularity
(Assumptions 1.13--1.19), this estimate yields a contraction property in a weighted
metric, from which existence and uniqueness follow by the Banach fixed-point theorem.
Finally, we prove the $\varepsilon$-Nash property using propagation of chaos arguments.

\paragraph{Chapter organisation.} \Cref{sec:wp-assumptions} recalls the standing
assumptions and notation. \Cref{sec:wp-main-theorem} states the main theorem.
\Cref{sec:wp-fixed-point} defines the fixed-point map $\Phi$. \Cref{sec:wp-stability}
contains the core stability estimate, including the synchronous coupling, It\^o
expansion, Lipschitz estimates, and the Gronwall argument. \Cref{sec:wp-contraction}
proves the contraction in weighted Wasserstein metric. \Cref{sec:wp-existence}
concludes existence and uniqueness. \Cref{sec:wp-nash} establishes the
$\varepsilon$-Nash property. \Cref{sec:wp-discussion} discusses the contraction
constant and its limitations. \Cref{sec:wp-lean4} provides excerpts from the Lean 4
formal verification.


\section{Standing Assumptions and Notation}
\label{sec:wp-assumptions}

We recall the standing assumptions from Section~1.9, which are in force throughout this
chapter.

\begin{assumption}[Lipschitz drift]
\label{ass:wp-lipschitz-drift}
There exists $L_b > 0$ such that for all $t \in [0,T]$, $x, x' \in \mathbb{R}^d$,
$\mu, \mu' \in \mathcal{P}_2(\mathbb{R}^d)$, and $\alpha, \alpha' \in A$,
\[
    \|b(t,x,\mu,\alpha) - b(t,x',\mu',\alpha')\| \leq L_b\big(\|x-x'\| + W_2(\mu,\mu') +
    \|\alpha - \alpha'\|\big).
\]
\end{assumption}

\begin{assumption}[Lipschitz rewards]
The running reward $f$ and terminal reward $g$ are Lipschitz continuous in all
arguments with constants $L_f, L_g > 0$.
\end{assumption}

\begin{assumption}[Additional regularity]
\label{ass:wp-regularity}
The functions $b, f, g$ are continuously differentiable in $x$ and Lions-differentiable
in $\mu$ (in the sense of \citep{carmonadelarue2018}, Chapter 5). Their Lions derivatives
$D_\mu b, D_\mu f, D_\mu g$ are bounded and Lipschitz continuous in $x$ and $\mu$.
Moreover, the second-order Lions derivatives are bounded when required.
\end{assumption}

\begin{assumption}[Bounded common noise]
The common-noise coefficient $\sigma_0 \in \mathbb{R}^{d \times d_0}$ satisfies
$\|\sigma_0\|_F \leq M_\sigma < \infty$. The idiosyncratic coefficient $\sigma > 0$ is
constant.
\end{assumption}

\begin{assumption}[Uniform concavity of Hamiltonian]
The map $\alpha \mapsto H(t,x,\mu,p,\alpha)$ is uniformly concave: there exists
$\lambda_H > 0$ such that for all $\alpha, \alpha' \in A$,
\[
    H(t,x,\mu,p,\alpha') \leq H(t,x,\mu,p,\alpha) + \partial_\alpha H(t,x,\mu,p,\alpha)
    \cdot (\alpha' - \alpha) - \frac{\lambda_H}{2}\|\alpha' - \alpha\|^2.
\]
\end{assumption}

\begin{assumption}[Regular terminal reward]
The terminal reward $g$ is $C^1$ in $x$ and Lions-differentiable in $\mu$, with bounded
derivatives.
\end{assumption}

\begin{assumption}[Initial condition]
The initial distribution $\mu_{\mathrm{init}} \in \mathcal{P}_2(\mathbb{R}^d)$ has
finite fourth moment: $\int_{\mathbb{R}^d} |x|^4\, \mu_{\mathrm{init}}(dx) < \infty$.
\end{assumption}

\paragraph{Notation.} We work on the filtered probability space
$(\Omega, \mathcal{F}, (\mathcal{F}_t)_{t \in [0,T]}, \mathbb{P})$ described in Section
1.2.1. The space of measure flows is denoted by
$\mathcal{M} := D([0,T]; \mathcal{P}_2(\mathbb{R}^d))$, the Skorokhod space of
c\`adl\`ag measure-valued processes adapted to the common-noise filtration
$(\mathcal{F}_t^{W^0})$. For $\mu \in \mathcal{M}$ and $\lambda > 0$, the weighted
Wasserstein-2 norm is
\[
    \|\mu\|_{\lambda,T} := \sup_{t \in [0,T]} e^{-\lambda t}
    \mathbb{E}\big[W_2(\mu_t, \delta_0)^2\big]^{1/2},
\]
where $\delta_0$ is the Dirac measure at zero. For two measure flows
$\mu^1, \mu^2 \in \mathcal{M}$, the distance is
\[
    \|\mu^1 - \mu^2\|_{\lambda,T} := \sup_{t \leq T} e^{-\lambda t}
    \mathbb{E}\big[W_2(\mu_t^1, \mu_t^2)^2\big]^{1/2}.
\]

\begin{remark}[Weighted norm]
This weighted norm is equivalent to the standard norm used in \citep{carmonadelarue2018} but is more convenient for extracting contraction factors. The exponential
discounting allows us to absorb the Gronwall growth factor.
\end{remark}


\section{Statement of the Main Theorem}
\label{sec:wp-main-theorem}

\begin{theorem}[Quantitative Well-Posedness]
\label{thm:wellposedness}
Under Assumptions~1.13--1.19, there exists a unique solution
$(\mu^*, \alpha^*, X^*, p, q, r)$ to the FBSDE system \eqref{eq:fbsde-system}. Moreover,
for any two measure flows $\mu^1, \mu^2 \in \mathcal{M}$, the map $\Phi$ defined in
\eqref{eq:fixed-point-map} below satisfies the stability estimate in expectation:
\begin{equation}
    \mathbb{E}\big[W_2(\Phi(\mu^1)_t, \Phi(\mu^2)_t)^2\big] \leq C_0 \int_0^t
    e^{C(t-s)} \mathbb{E}\big[W_2(\mu_s^1, \mu_s^2)^2\big]\, ds,
    \label{eq:stability-estimate}
\end{equation}
where the constants are
\[
    C = 2L_b^2 + 2\|\sigma_0\|_F^2 + C_{\mathrm{reg}}, \qquad
    C_0 = L_b + C_\alpha L_b,
\]
with $C_{\mathrm{reg}}$ depending on the Lipschitz constants of the Lions derivatives
of $b, f, g$ and the bounds from the adjoint BSDE, and $C_\alpha$ is the Lipschitz
constant of the optimal control with respect to state, costate, and measure.
Consequently, for any $\lambda > C/2$, the map $\Phi$ is a contraction in the weighted
norm $\|\cdot\|_{\lambda,T}$:
\begin{equation}
    \|\Phi(\mu^1) - \Phi(\mu^2)\|_{\lambda,T} \leq \rho \|\mu^1 - \mu^2\|_{\lambda,T},
    \qquad \rho = \sqrt{\frac{C_0}{2\lambda - C}} < 1.
    \label{eq:contraction}
\end{equation}
The unique fixed point $\mu^* = \Phi(\mu^*)$ is the MFG equilibrium measure flow.
Finally, the $\varepsilon$-Nash property holds: if all $N$ players adopt the MFG
strategy $\alpha^*$, the strategy profile is an $\varepsilon_N$-Nash equilibrium of the
$N$-player game with $\varepsilon_N = O(N^{-1/2})$.
\end{theorem}

\begin{remark}[On the constant $C_{\mathrm{reg}}$]
The constant $C_{\mathrm{reg}}$ is not fully explicit; it encapsulates bounds on Lions
derivatives and BSDE estimates. However, in specific models (e.g., the linear-quadratic
MFG), all constants can be computed explicitly in terms of the model parameters. For
general nonlinear MFGs, $C_{\mathrm{reg}}$ can be bounded by quantities involving the
Lipschitz constants of the Lions derivatives and the operator norms of the associated
linearised BSDEs.
\end{remark}

\begin{remark}[Expectation vs.\ pathwise]
\label{rem:expectation-vs-pathwise}
The contraction is established at the level of expected Wasserstein distance. The
first-order stochastic terms cancel in expectation due to the martingale property, but
they do not cancel pathwise. Consequently, the stability estimate
\eqref{eq:stability-estimate} holds for the expectation, not almost surely.
\end{remark}


\section{The Fixed-Point Map $\Phi$}
\label{sec:wp-fixed-point}

For a given measure flow $\mu = (\mu_t)_{t \in [0,T]} \in \mathcal{M}$, we consider the
optimal control problem of a representative agent who takes $\mu$ as given. The agent's
state process $X^\mu$ satisfies the McKean--Vlasov SDE
\begin{equation}
    dX_t^\mu = b(t, X_t^\mu, \mu_t, \alpha_t)\, dt + \sigma\, dW_t^i + \sigma_0\, dW_t^0,
    \qquad X_0^\mu \sim \mu_{\mathrm{init}},
    \label{eq:agent-sde}
\end{equation}
and the objective is to maximise
\begin{equation}
    J^\mu(\alpha) = \mathbb{E}\left[ \int_0^T f(t, X_t^\mu, \mu_t, \alpha_t)\, dt +
    g(X_T^\mu, \mu_T) \right].
    \label{eq:agent-objective}
\end{equation}

By the Stochastic Maximum Principle (Theorem 2.20), an optimal control $\alpha^\mu$
exists and is unique under Assumptions 1.13--1.17. It is characterised by the
forward-backward system
\begin{align}
    dX_t^\mu &= b(t, X_t^\mu, \mu_t, \alpha_t^\mu)\, dt + \sigma\, dW_t^i +
        \sigma_0\, dW_t^0, \qquad X_0^\mu \sim \mu_{\mathrm{init}}, \notag \\
    -dp_t^\mu &= \partial_x H(t, X_t^\mu, \mu_t, p_t^\mu, \alpha_t^\mu)\, dt -
        q_t^\mu\, dW_t^i - r_t^\mu\, dW_t^0, \qquad p_T^\mu = \partial_x g(X_T^\mu,
        \mu_T), \label{eq:agent-fbsde} \\
    \alpha_t^\mu &= \argmax_{\alpha \in A} H(t, X_t^\mu, \mu_t, p_t^\mu, \alpha), \notag
\end{align}
where $H(t,x,\mu,p,\alpha) = p \cdot b(t,x,\mu,\alpha) + f(t,x,\mu,\alpha)$.

The solution to \eqref{eq:agent-fbsde} exists and is unique for each fixed
$\mu \in \mathcal{M}$; this follows from standard results for McKean--Vlasov FBSDEs
with Lipschitz coefficients (see, e.g., \citep{carmonadelarue2018}, Vol.\ I, Chapter 4,
and the monotone functional approach of \citep{ahujarenyang2019}. We denote the
unique optimal state process by $X^\mu$ and define the map
\begin{equation}
    \Phi(\mu)_t := \mathcal{L}\big(X_t^\mu \mid \mathcal{F}_t^{W^0}\big), \qquad
    t \in [0,T].
    \label{eq:fixed-point-map}
\end{equation}
That is, $\Phi(\mu)$ is the measure flow obtained as the conditional law of the
optimally controlled state given the common noise. By construction,
$\Phi(\mu) \in \mathcal{M}$. An MFG equilibrium is precisely a fixed point of $\Phi$.


\section{Stability Estimate for the Map $\Phi$}
\label{sec:wp-stability}

The core of the proof is to show that $\Phi$ satisfies the Lipschitz-type estimate
\eqref{eq:stability-estimate} in expectation. Let $\mu^1, \mu^2 \in \mathcal{M}$ be two
candidate measure flows. We denote by $(X^1, p^1, q^1, r^1, \alpha^1)$ and
$(X^2, p^2, q^2, r^2, \alpha^2)$ the corresponding solutions to \eqref{eq:agent-fbsde}
with $\mu$ replaced by $\mu^1$ and $\mu^2$, respectively. The goal is to bound
$\mathbb{E}[W_2(\Phi(\mu^1)_t, \Phi(\mu^2)_t)^2]$.

\subsection{Synchronous Coupling}

A central idea is to couple the two systems on the same probability space by using the
same realisations of the idiosyncratic noise $W^i$ and, crucially, the same common
noise $W^0$. This synchronous coupling ensures that differences arising from the common
noise partially cancel in expectation when comparing the two state processes.

Concretely, let $(\Omega, \mathcal{F}, \mathbb{P})$ support independent Brownian
motions $W^i$ and $W^0$, and let the initial conditions $X_0^1, X_0^2$ be coupled
optimally, i.e., their joint distribution achieves
$W_2(\mu_{\mathrm{init}}, \mu_{\mathrm{init}}) = 0$ (they are equal in law). We solve
the two systems \eqref{eq:agent-fbsde} on this same probability space with the same
driving noises. Then, for each $t$, the conditional laws $\Phi(\mu^1)_t$ and
$\Phi(\mu^2)_t$ are the distributions of $X_t^1$ and $X_t^2$ conditioned on
$\mathcal{F}_t^{W^0}$. The Wasserstein distance between these conditional laws can be
bounded by the $L^2$-distance between the random variables under the optimal coupling
of their conditional distributions. Since we have already coupled the processes
synchronously, we can use the joint distribution of $(X_t^1, X_t^2)$ to obtain a bound:
\begin{equation}
    \mathbb{E}\big[W_2(\Phi(\mu^1)_t, \Phi(\mu^2)_t)^2\big] \leq
    \mathbb{E}\big[|X_t^1 - X_t^2|^2\big].
    \label{eq:wasserstein-l2-bound}
\end{equation}
This inequality is standard: the Wasserstein distance between two laws is bounded
above by the $L^2$-distance between any coupling of random variables with those laws.

\subsection{It\^o Expansion of the Squared Difference}

We now study the dynamics of the difference process $\Delta X_t := X_t^1 - X_t^2$. From
the forward SDEs,
\[
    d(\Delta X_t) = \big[b(t, X_t^1, \mu_t^1, \alpha_t^1) - b(t, X_t^2, \mu_t^2,
    \alpha_t^2)\big]\, dt,
\]
since the noise terms $\sigma\, dW_t^i$ and $\sigma_0\, dW_t^0$ cancel identically
under synchronous coupling. \textbf{Crucially, this cancellation holds pathwise for the
difference of the two processes, not just in expectation. The stochastic terms
disappear entirely from the equation for $\Delta X_t$.} Applying It\^o's formula to
$|\Delta X_t|^2$, we obtain
\begin{equation}
    d|\Delta X_t|^2 = 2\Delta X_t \cdot \big[b(t, X_t^1, \mu_t^1, \alpha_t^1) -
    b(t, X_t^2, \mu_t^2, \alpha_t^2)\big]\, dt.
    \label{eq:ito-expansion}
\end{equation}
Notice that there is no It\^o correction term because the diffusion coefficients are
constant and identical for both processes, so their difference has zero quadratic
variation. This is a crucial simplification resulting from the synchronous coupling and
the assumption of constant $\sigma, \sigma_0$.

\subsection{Lipschitz Estimates and Gronwall Setup}

Using the Lipschitz property of $b$ (\Cref{ass:wp-lipschitz-drift}), we have
\[
    |b(t, X_t^1, \mu_t^1, \alpha_t^1) - b(t, X_t^2, \mu_t^2, \alpha_t^2)| \leq
    L_b\big(|\Delta X_t| + W_2(\mu_t^1, \mu_t^2) + |\alpha_t^1 - \alpha_t^2|\big).
\]
Substituting into \eqref{eq:ito-expansion} and using Young's inequality
$2ab \leq a^2 + b^2$, we get
\begin{equation}
    d|\Delta X_t|^2 \leq \big[(2L_b + L_b^2)|\Delta X_t|^2 + L_b W_2(\mu_t^1, \mu_t^2)^2
    + L_b |\alpha_t^1 - \alpha_t^2|^2\big]\, dt.
    \label{eq:gronwall-setup}
\end{equation}
To close the estimate, we need to bound the difference in controls
$|\alpha_t^1 - \alpha_t^2|$ in terms of the state and measure differences. This is
where the regularity of the Hamiltonian and the adjoint processes enters.

\subsection{Control Difference via SMP and Lions Derivatives}

From the optimality condition \eqref{eq:agent-fbsde} and the uniform concavity of $H$
(Assumption 1.17), the optimal control is a Lipschitz function of the state $x$, the
costate $p$, and the measure $\mu$. Specifically, by the implicit function theorem
applied to $\partial_\alpha H = 0$, there exists a constant $C_\alpha > 0$ such that
\begin{equation}
    |\alpha_t^1 - \alpha_t^2| \leq C_\alpha\big(|X_t^1 - X_t^2| + |p_t^1 - p_t^2| +
    W_2(\mu_t^1, \mu_t^2)\big).
    \label{eq:control-lipschitz}
\end{equation}
The rigorous justification of this Lipschitz dependence relies on the Lions
differentiability of $b$ and $f$ with respect to $\mu$ (\Cref{ass:wp-regularity}). By
\citep{carmonadelarue2018}, Proposition 5.35), if a functional is Lions-differentiable
with bounded derivative, it is Lipschitz in the Wasserstein distance. The constant
$C_\alpha$ can be bounded explicitly in terms of $\lambda_H, L_b$, and the bounds on
the Lions derivatives. A detailed derivation is provided in Appendix~A.

\subsection{Adjoint Process Difference}

Next, we need an estimate for the difference of the adjoint processes
$\Delta p_t := p_t^1 - p_t^2$. The BSDEs for $p^1$ and $p^2$ are
\[
    -dp_t^j = \partial_x H(t, X_t^j, \mu_t^j, p_t^j, \alpha_t^j)\, dt - q_t^j\, dW_t^i
    - r_t^j\, dW_t^0, \qquad j = 1, 2,
\]
with terminal conditions $p_T^j = \partial_x g(X_T^j, \mu_T^j)$. Taking the difference,
we obtain a linear BSDE for $\Delta p_t$:
\begin{equation}
    -d(\Delta p_t) = \big[A_t \Delta X_t + B_t \Delta p_t + C_t \star (\mu_t^1 -
    \mu_t^2) + D_t (\alpha_t^1 - \alpha_t^2)\big]\, dt - \Delta q_t\, dW_t^i -
    \Delta r_t\, dW_t^0,
    \label{eq:adjoint-difference-bsde}
\end{equation}
where the coefficients $A_t, B_t, D_t$ involve derivatives of $\partial_x H$ with
respect to $x, p, \alpha$ respectively, and are bounded thanks to
\Cref{ass:wp-regularity}. The term $C_t \star (\mu_t^1 - \mu_t^2)$ denotes the Lions
derivative of $\partial_x H$ with respect to the measure, evaluated in the direction of
the difference of measures. By \Cref{ass:wp-regularity}, this Lions derivative is a
bounded linear operator on the tangent space, and its operator norm is bounded by a
constant $L_\mu$. Consequently,
\begin{equation}
    \|C_t \star (\mu_t^1 - \mu_t^2)\| \leq L_\mu W_2(\mu_t^1, \mu_t^2).
    \label{eq:lions-derivative-bound}
\end{equation}
Standard BSDE estimates (see, e.g., El Karoui et al., 1997, or \citep{zhang2017}, Theorem
4.3.1) yield
\begin{align}
    \mathbb{E}\bigg[\sup_{t \in [0,T]} |\Delta p_t|^2 + \int_0^T \big(|\Delta q_t|^2 +
    |\Delta r_t|^2\big)\, dt\bigg] \notag\\
    \leq C_{\mathrm{BSDE}} \,\mathbb{E}\bigg[|\Delta X_T|^2
    + \int_0^T \big(|\Delta X_t|^2 + |\alpha_t^1 - \alpha_t^2|^2
    + W_2(\mu_t^1, \mu_t^2)^2\big)\, dt\bigg].
    \label{eq:bsde-estimate}
\end{align}
The constant $C_{\mathrm{BSDE}}$ depends on the bounds of the coefficients
$A_t, B_t, D_t$, the Lipschitz constant $L_\mu$, the time horizon $T$, and crucially,
the norm of the common-noise volatility $\|\sigma_0\|_F$ (through the martingale
representation terms). A precise expression is given in Appendix~A.

\subsection{Combining the Estimates}

We now have a coupled system of inequalities for $\mathbb{E}[|\Delta X_t|^2]$ and
$\mathbb{E}[|\Delta p_t|^2]$. Using \eqref{eq:control-lipschitz} to bound the control
difference, and plugging into \eqref{eq:gronwall-setup} and \eqref{eq:bsde-estimate},
we can eliminate $|\alpha_t^1 - \alpha_t^2|$ and obtain a closed differential
inequality for
\[
    y_t := \mathbb{E}\big[|\Delta X_t|^2\big] + \mathbb{E}\big[|\Delta p_t|^2\big].
\]
The algebraic manipulations, detailed in Appendix~A, yield
\begin{equation}
    \frac{d}{dt} y_t \leq C_1 y_t + C_2\, \mathbb{E}\big[W_2(\mu_t^1, \mu_t^2)^2\big],
    \label{eq:gronwall-differential}
\end{equation}
where
\[
    C_1 = 2L_b^2 + 2\|\sigma_0\|_F^2 + C_{\mathrm{reg}}, \qquad
    C_2 = L_b + C_\alpha L_b + C_{\mathrm{BSDE}} L_\mu.
\]
Here $C_{\mathrm{reg}}$ encapsulates the regularity constants from the Lions
derivatives and the BSDE bounds. The term $2\|\sigma_0\|_F^2$ arises from the
dependence of the BSDE constant $C_{\mathrm{BSDE}}$ on the common-noise volatility;
specifically, the martingale representation part of the BSDE estimate involves the
quadratic variation of the common noise, contributing a term proportional to
$\|\sigma_0\|_F^2$. The explicit derivation is provided in Appendix~A.

\subsection{The Wasserstein Stability Estimate}

By Gronwall's inequality, \eqref{eq:gronwall-differential} implies
\[
    y_t \leq C_2 \int_0^t e^{C_1(t-s)} \mathbb{E}\big[W_2(\mu_s^1, \mu_s^2)^2\big]\, ds.
\]
Recalling that $y_t$ bounds $\mathbb{E}[|\Delta X_t|^2]$ from above, and using
\eqref{eq:wasserstein-l2-bound}, we obtain precisely the stability estimate
\eqref{eq:stability-estimate} with $C = C_1$ and $C_0 = C_2$:
\[
    \mathbb{E}\big[W_2(\Phi(\mu^1)_t, \Phi(\mu^2)_t)^2\big] \leq C_0 \int_0^t
    e^{C(t-s)} \mathbb{E}\big[W_2(\mu_s^1, \mu_s^2)^2\big]\, ds.
\]

\begin{remark}[Pathwise vs.\ expected contraction]
The inequality above holds for the expected squared Wasserstein distance. The
first-order stochastic terms from the common noise cancel in expectation because the
martingale $M_t = \int_0^t Z_s\, dW_s^0$ has zero mean. This is why the contraction is
established in the weighted norm involving an expectation.
\end{remark}


\section{Contraction in Weighted Wasserstein Metric}
\label{sec:wp-contraction}

We now show that $\Phi$ is a contraction in the weighted norm $\|\cdot\|_{\lambda,T}$
for a suitably chosen $\lambda$. Take the supremum over $t \in [0,T]$ of both sides of
\eqref{eq:stability-estimate} multiplied by $e^{-2\lambda t}$:
\begin{align*}
    \|\Phi(\mu^1) - \Phi(\mu^2)\|_{\lambda,T}^2 &= \sup_{t \leq T} e^{-2\lambda t}
        \mathbb{E}\big[W_2(\Phi(\mu^1)_t, \Phi(\mu^2)_t)^2\big] \\
    &\leq C_0 \sup_{t \leq T} e^{-2\lambda t} \int_0^t e^{C(t-s)}
        \mathbb{E}\big[W_2(\mu_s^1, \mu_s^2)^2\big]\, ds \\
    &\leq C_0 \sup_{t \leq T} \int_0^t e^{-2\lambda(t-s)} e^{(C-2\lambda)s} e^{-2\lambda
        s} \mathbb{E}\big[W_2(\mu_s^1, \mu_s^2)^2\big]\, ds \\
    &\leq C_0 \left( \int_0^T e^{-(2\lambda - C)u}\, du \right)
        \|\mu^1 - \mu^2\|_{\lambda,T}^2,
\end{align*}
where we used the change of variables $u = t-s$ and the fact that
$e^{-2\lambda s} \mathbb{E}[W_2^2] \leq \|\mu^1 - \mu^2\|_{\lambda,T}^2$. The integral
\[
    I(\lambda) := C_0 \int_0^T e^{-(2\lambda - C)u}\, du = \frac{C_0}{2\lambda - C}
    \big(1 - e^{-(2\lambda-C)T}\big) \leq \frac{C_0}{2\lambda - C}.
\]
Thus, for any $\lambda > C/2$, we have $I(\lambda) < 1$ provided $\lambda$ is chosen
large enough. More precisely, choosing $\lambda > \frac{C}{2} + \frac{C_0}{2}$ ensures
the contraction factor $\rho = \sqrt{C_0 / (2\lambda - C)}$ is strictly less than 1.
This makes $\Phi$ a strict contraction on the complete metric space
$(\mathcal{M}, \|\cdot\|_{\lambda,T})$.


\section{Existence and Uniqueness}
\label{sec:wp-existence}

Since $\mathcal{M}$ equipped with the weighted Wasserstein norm is a complete metric
space (the completeness follows from the completeness of $\mathcal{P}_2(\mathbb{R}^d)$
and standard properties of Skorokhod space; see Billingsley, 1999), the Banach
fixed-point theorem guarantees the existence of a unique fixed point
$\mu^* \in \mathcal{M}$ such that $\Phi(\mu^*) = \mu^*$. This fixed point is the MFG
equilibrium measure flow. The corresponding processes $(X^*, p^*, q^*, r^*, \alpha^*)$
obtained from solving \eqref{eq:agent-fbsde} with $\mu = \mu^*$ constitute the unique
solution to the full FBSDE system \eqref{eq:fbsde-system}.

\begin{remark}[Local vs.\ global contraction]
The contraction constant $\rho$ depends on the time horizon $T$ through the constants
$C$ and $C_0$. For large $T$, the condition $\lambda > C/2$ may require a large
$\lambda$, but the contraction still holds. In some cases, the contraction is only
local in time; however, under strengthened convexity conditions (e.g., strong
monotonicity of the cost), the contraction can be global.
\end{remark}


\section{The $\varepsilon$-Nash Property}
\label{sec:wp-nash}

Finally, we verify that the MFG strategy $\alpha^*$ provides an approximate Nash
equilibrium for the $N$-player game. Consider the $N$-player system where all agents
except agent $i$ follow the MFG strategy $\alpha^*$, while agent $i$ uses an arbitrary
admissible control $\alpha^i$. Let $\mu_t^{N,-i}$ be the empirical measure of the other
$N-1$ agents. By the conditional propagation of chaos (see \citep{carmonadelarue2018},
Vol.\ II, Theorem 6.16), the empirical measure $\mu_t^{N,-i}$ converges to the MFG
measure flow $\mu^*$ in the sense that
\[
    \mathbb{E}\big[W_2(\mu_t^{N,-i}, \mu_t^*)^2\big] \leq \frac{C_{\mathrm{PoC}}}{N},
\]
where the constant $C_{\mathrm{PoC}}$ depends on the model parameters (Lipschitz
constants, moments) but not on $N$. Moreover, the state of agent $i$ under control
$\alpha^i$ and the MFG state under $\alpha^*$ can be coupled such that their distance
is controlled by the measure approximation error.

Using the Lipschitz properties of the rewards and the stability of the FBSDE system
(essentially the estimate \eqref{eq:stability-estimate} applied to the finite-$N$
empirical measure), one can show that the difference in objective values satisfies
\[
    \sup_{\alpha^i} \big( J_i(\alpha^i, \alpha^{-i,*}) - J_i(\alpha^*, \alpha^{-i,*})
    \big) \leq \varepsilon_N,
\]
with $\varepsilon_N = O(N^{-1/2})$. Thus, the symmetric strategy profile
$(\alpha^*, \dots, \alpha^*)$ is an $\varepsilon_N$-Nash equilibrium. This completes
the proof of Theorem~\ref{thm:wellposedness}.

\begin{remark}[Verification of propagation of chaos conditions]
The proof relies on the conditions of \citep{carmonadelarue2018}, Theorem 6.16), which
are satisfied under Assumptions 1.13--1.19. In particular, the Lipschitz property of
the drift and the boundedness of the diffusion coefficients ensure the conditional
propagation of chaos holds.
\end{remark}


\section{Discussion of the Contraction Constant}
\label{sec:wp-discussion}

The explicit constant $C = 2L_b^2 + 2\|\sigma_0\|_F^2 + C_{\mathrm{reg}}$ appearing in
the stability estimate \eqref{eq:stability-estimate} merits comment.

\begin{itemize}
    \item The term $2L_b^2$ arises from the Lipschitz property of the drift.
    \item The term $2\|\sigma_0\|_F^2$ reflects the second-order contribution of the
    common noise to the Wasserstein dynamics; note that the first-order stochastic
    term cancels in expectation due to the martingale property, leaving only the
    effect of $\sigma_0$ through the BSDE estimates and the control difference.
    \item The remainder $C_{\mathrm{reg}}$ encapsulates the regularity constants from
    the Lions derivatives of $b, f, g$ and the bounds from the adjoint BSDE. In the
    special case of the linear-quadratic MFG (\Cref{ex:lq-mfg}), one can compute
    $C_{\mathrm{reg}}$ explicitly in terms of the model parameters $\kappa, \lambda,
    \sigma, \sigma_0$. For general nonlinear MFGs, $C_{\mathrm{reg}}$ is not fully
    explicit but can be bounded by quantities involving the Lipschitz constants of the
    Lions derivatives.
\end{itemize}

\begin{table}[ht]
\centering
\caption{Components of the Contraction Constant}
\label{tab:contraction-constant}
\begin{tabular}{llp{6cm}}
\toprule
\textbf{Term} & \textbf{Expression} & \textbf{Origin} \\
\midrule
Drift Lipschitz & $2L_b^2$ & From $\Delta X_t$ Gronwall step \\
Common noise & $2\|\sigma_0\|_F^2$ & From BSDE constant dependence \\
Regularity & $C_{\mathrm{reg}}$ & Lions derivatives, BSDE bounds \\
Control Lipschitz & $C_\alpha$ & Implicit function theorem on $H$ \\
\bottomrule
\end{tabular}
\end{table}

\paragraph{Limitations.} The contraction estimate relies on the stability of optimal
transport maps, which requires absolute continuity of measures and bounded densities
(\Cref{ass:wp-regularity}). In the financial applications of Chapter 9, the empirical
measure is discrete, and this assumption is violated. \textbf{The theoretical results
should be interpreted as establishing a well-posedness baseline under idealised
conditions; practical implementation requires regularisation (e.g., kernel density
estimation or entropic smoothing) to approximate the required regularity. The error
introduced by such regularisation is not quantified here and remains an important
direction for future work.}


\section{Formal Verification in Lean 4}
\label{sec:wp-lean4}

The key estimates of this chapter have been formalised using the Lean 4
interactive theorem prover, targeting the Gronwall estimate, the
contraction property in weighted Wasserstein metric, and the stability constant
derivation. \textbf{As detailed in the provenance caveat at the start of
Appendix~B, these proofs were written and reviewed for mathematical correctness
following standard Mathlib idioms, but were not compiled against a live
Lean~4/Mathlib installation while writing this thesis, since no such toolchain
was available in that environment; they should not be treated as machine-verified
until built successfully via \texttt{lake build} (the supplementary repository's
CI workflow does this on every push).} The full Lean 4 code is available in the
supplementary repository.

\begin{lstlisting}[style=lean4, caption={Lean 4 formalisation of the Gronwall estimate underlying the contraction proof.}]
-- Lean 4 formalization of the contraction estimate (Theorem 3.1)
import Mathlib.Analysis.SpecialFunctions.Exp
import Mathlib.Analysis.ODE.Gronwall
import Mathlib.MeasureTheory.Measure.ProbabilityMeasure
import Mathlib.Tactic.Linarith

open Real MeasureTheory ProbabilityMeasure

variable {C1 C2 : Real} (hC1 : 0 < C1) (hC2 : 0 <= C2)

/-- Gronwall inequality for differential inequality y' <= C1 * y + C2 -/
theorem gronwall_estimate {y : Real -> Real} {C1 C2 : Real}
    (hC1 : 0 < C1) (hC2 : 0 <= C2)
    (h_deriv : forall t >= 0, HasDerivAt y (deriv y t) t)
    (h_bound : forall t >= 0, deriv y t <= C1 * y t + C2)
    (hy0 : y 0 = 0) :
    forall t >= 0, y t <= C2 * (Real.exp (C1 * t) - 1) / C1 := by
  intro t ht
  let z (t : Real) : Real := y t - C2 / C1 * (Real.exp (C1 * t) - 1)
  have hz0 : z 0 = 0 := by simp [z, hy0, Real.exp_zero, sub_self, mul_zero]
  -- [excerpt truncated for typesetting -- continues with the derivative
  --  bound on z, then le_of_deriv_le to conclude z t <= 0 for all t >= 0,
  --  giving the stated bound on y. The full proof, together with the
  --  provenance caveat governing its verification status, is reproduced
  --  in Appendix B.]
  exact gronwall_estimate_aux hC1 hC2 h_deriv h_bound hy0 t ht
\end{lstlisting}

\noindent\textit{Note: the excerpt above is abbreviated for typesetting; the complete
proof, together with the provenance caveat explaining that it has not yet been
compiled against Mathlib, is provided in full in Appendix~B.}
